% /chapters/2-descriptive-combinatorics/02-concise-L0.tex

\section{A concise proof of the \texorpdfstring{$L_0$}{L0} dichotomy}\label{sec:concise_L0}

Carroy, Miller, Schrittesser, and Vidny\'anszky established the $L_0$ dichotomy: there is a Borel graph of Borel chromatic number three that admits a continuous homomorphism to every analytic graph of Borel chromatic number at least three.
Their proof relies on a transfinite analysis of terminal approximations over a decreasing $\omega_1$-sequence of analytic sets.

In this section, I give a new, substantially shorter proof of this result by adapting the graph-theoretic framework recently introduced by Bernshteyn for the $G_0$ dichotomy.
The central device is a $\sigma$-ideal of \emph{small} sets of homomorphisms from finite path approximations into the target graph, where smallness is witnessed by a bounded odd-walk condition on vertex projections.
The key lemma that largeness is preserved under the doubling operation is established via the First Reflection Theorem, replacing the original transfinite construction with a single Borel reflection argument.
The continuous homomorphism from the canonical graph $\mathbb L_c$ into the target is then obtained as a limit of shrinking families of copies, in direct analogy with Bernshteyn's proof for $G_0$.

\subsection{Introduction}\label{sec:concise-intro}

Throughout this entire section, $G$ denotes an analytic graph on a Polish space $V(G)=X$.
That is, $X$ is a separable completely metrizable topological space, and $G$ is a symmetric and irreflexive relation on $X$ which, as a subspace of $X\times X$, is the continuous image of a Polish space.

Denote the \emph{Borel chromatic number} of $G$ by $\chi_B(G)$, that is, $\chi_B(G)$ is the least number of colours $k$ needed to find a Borel mapping $c\colon X\to k$ in such a way that adjacent vertices in $G$ are given different colours.
I employ the symbol $H\to_cG$ (respectively, $H\to_BG$) to abbreviate the fact that there is a continuous (respectively Borel) homomorphism from the analytic graph $H$ to the analytic graph $G$.
It is straightforward to show the following.

\begin{fact}\label{fact:chrom_transitive}
    If $H\to_cG$ or, more generally if $H\to_BG$, then $\chi_B(H)\leq\chi_B(G)$.
\end{fact}

Using a standard greedy algorithm argument, it is clear that any finite graph $G$ of maximum degree $\Delta$ satisfies $\chi(G)\leq\Delta+1$.
Interestingly, a similar fact holds for Borel graphs of uniformly bounded degree.

\begin{theorem}[Proposition~4.6 in \cite{kechris1999borel}]\label{prop:brooks}
    If $G$ is a Borel graph on a Polish space all of whose degrees are bounded by the natural number $k$, then $\chi_B(G)\leq k+1$.
\end{theorem}

The following notion, taken from \cite[Theorem~35.10, p.~285]{Kechris_1995_DescriptiveSetTheory}, is a well-known fact about analytic sets.

\begin{definition}\label{def:pi11-on-sigma11}
    A collection of subsets $\Phi$ of a Polish space $X$ is said to be \emph{$\Pi^1_1$ on $\Sigma^1_1$} if for any Polish $Y$ and any $\Sigma^1_1$ set $A\subseteq Y\times X$, the set $A_\Phi:=\{y\in Y:A_y:=\{x\in X:(y,x)\in A\}\in\Phi\}$ is $\Pi^1_1$.
\end{definition}

\begin{lemma}[First Reflection Theorem]\label{lemma:first-reflection}
    If $\Phi$ is as in the preceding definition, then for any $\Sigma^1_1$ set $A\in\Phi$, there is a Borel set $B\subseteq X$ such that $A\subseteq B$ and $B\in\Phi$.
\end{lemma}

\subsection{Finite path approximations}\label{sec:finite-path-approx}

Fix the parameter $c\in\omega^\omega$ and a path of length $\omega$ on the vertex set $\{p_n\}_{n<\omega}$.
That is, the edge set $\{(p_n,p_{n+1}):n<\omega\}$.
Naturally, one could assume that $p_n=n$, but the labeling helps distinguish these vertices from other indices.

Recall that if $s$ and $t$ are finite strings (of, say, natural numbers), then $s^\frown t$ is the \emph{concatenation} of $s$ and $t$.
Similarly, $t^n$ is the repeated concatenation of $t$ with itself, and $(i)$ denotes the string of single value $i$.
Next, recursively construct a family of graphs $\{L^c_n:n<\omega\}$ such that for all $n<\omega$:
$L^c_n$ is a finite path; if $n>0$, the endpoints of $L^c_n$ are $e_i^n:=(p_0)^\smallfrown(0)^{n-1}{}^\smallfrown(i)$, for $i<2$; and
\begin{equation}
\text{Every vertex of }L^c_n\text{ is of the form }(p_k)^\smallfrown t\text{ for some }k\text{ and }t\in2^{\leq n}.\label{eq:L_n-condition-iii}\end{equation}
Unless $t=\emptyset$, call such vertex a \emph{non-path vertex}, and call the vertices of the form $(p_k)$ the \emph{path vertices} of $L^c_n$.


Start with $L^c_0$ as the graph consisting of a single vertex $p_0$ and no edges, and set $e^0_0=e^0_1=p_0$.
Now suppose that $L^c_n$ has been constructed.
$L^c_{n+1}$ will be a path obtained from joining two copies of $L^c_n$ by a path of $c(n)+2$ edges (with $c(n)+1$ \emph{new} path vertices), but one must update the vertex labels in order to distinguish the two copies.
One accomplishes this by simply appending $0$ to each vertex of the first copy, and $1$ to the other.
Formally, start by adding the vertex $v^\frown(i)$ for every $i<2$ and $v\in L^c_n$.
Adding the path $$((e_1^n)^\smallfrown(0),p_0,p_1,\dots,p_{c(n)},(e_1^n)^\smallfrown(1))$$ completes the construction of $L^c_{n+1}$, and its endpoints are $$e_i^{n+1}=(e_0^n)^\smallfrown(i)=(p_0)^\smallfrown(0)^n{}^\smallfrown(i),$$ for $i<2$.
This finishes the recursion.
Note that to construct $L^c_{n+1}$ from $L^c_n$, one only requires knowledge of the value $c(n)$.

\begin{figure}[ht]
    \centering
    \input{figures/L0_figure.tex}
    \caption{$L_3^c$ for $c(0)=1,c(1)=3,c(2)=5$.}
    \label{fig:L_n}
\end{figure}

To develop some intuition on what \eqref{eq:L_n-condition-iii} says about a vertex $v$ in $L^c_n$: the number $|v|-1$ indicates that $v$ was added that many stages ago in the recursion, at the position of $p_k$ in the path joining the two copies of the previous paths; the sequence $t$ tells the story, in order, of which of the two copies of the previous path the vertex lies in.
In other words, vertices of the form $(p_k)^\smallfrown t$ for $t\in 2^n$ are precisely those that were copied from the original $L^c_0$ in all possible positions along $L^c_n$, doubling the number of copies at each stage.

Based on this discussion, a vertex in some $L^c_n$ is completely determined by the stage $n-m$ at which it was added in the position of $p_k$ and then copied according to the binary sequence $t$.

\subsection{A notion of smallness for copies}\label{sec:notion-smallness}

Now, define a notion of smallness analogous to the one presented in \cite{Bernshteyn2024}.
By definition, the set of edges of $G$ is the continuous image of some Polish space $E$, let us say under $\pi$.
For a finite graph $H$, consider $\Hom(H,G)$ as the set of all maps

\begin{align*}
    &\varphi\colon V(H)\to X\\
    &\varphi\colon E(H)\to E
\end{align*}

that satisfy that for all $uv\in E(H)$, $\pi(\varphi(uv))=\varphi(u)\varphi(v)$.\footnote{Formally, $\varphi$ is a pair of two maps, however one can assume that $V(H)\cap E(H)=\emptyset$ and then there is no risk of confusion.}
Recall that $X=V(G)$ is the Polish space of vertices on which the analytic graph $G$ is defined.
Notice that $\Hom(H,G)$ is also Polish space, being a closed subset of the product space $X^{V(H)}\times E^{E(H)}$.

Given $\mathcal H\subseteq\Hom(H,G)$, define $\mathcal H(u):=\{\varphi(u):\varphi\in\mathcal H\}$, which is analytic whenever $\mathcal H$ is Borel.
Also, let $\mathcal H(uv):=\{\varphi(uv):\varphi\in\mathcal H\}$.

\begin{definition}\label{def:small}~ 
\begin{enumerate}
    \item Let $k<\omega$.
    Say that the set $A$ has property $\Phi(A,k)$ if $A\subseteq X$ is an analytic set of vertices of $G$ such that all odd walks of $G$ with endpoints in $A$ have length at most $2k-1$.

    \item A Borel set $\mathcal L\subseteq\Hom(L^c_n,G)$ is called \emph{tiny} if there is a vertex $u\in V(L^c_n)$ and a natural $k$ such that $\Phi(\mathcal L(u),k)$.

    \item A set $\mathcal L\subseteq\Hom(L^c_n,G)$ is \emph{small} if it is the countable union of tiny sets.
    This notion defines a $\sigma$-ideal.

    \item A set is \emph{large} if it is not small.
    Notice that when $H$ is the graph with a single vertex $H=\bullet\left(=L^c_0\right)$, $\Hom(H,G)$ can be identified with $X$ and then for any $\mathcal H\subseteq\Hom(\bullet,G)$, $\mathcal H(\bullet)$ is identified with $\mathcal H$.
\end{enumerate}
\end{definition}

Now invoke the following auxiliary result that links the $\Phi$ property to Borel 2-colourability.
The proof proceeds by induction on $n$ using analytic separation; see \cite[Claims~3.3--3.5]{L0paper} for details.
Recall that a set of vertices of $G$ is \emph{invariant} if it is the union of connected components of $G$.
In addition, if $A\subseteq X$, then $[A]_{E_G}$ is the union of all connected components of $G$ that intersect $A$, that is, the set of vertices that can be reached from some element of $A$ by a finite path.
Therefore, a set $A$ is invariant if and only if $A=[A]_{E_G}$.

\begin{lemma}[Claim~3.5 of \cite{L0paper}]\label{lemma:L0-claim3.5}
    If $\Phi(A,n)$, then there is an invariant Borel set $B\supseteq[A]_{E_G}$ that induces a 2-Borel-colourable subgraph of $G$, that is, there is a Borel proper colouring $c_B\colon G\restriction_B\to2$.
\end{lemma}

\begin{proposition}\label{prop:small}
    If $\chi_B(G)>2$ then $X$ (via its identification with $\Hom(\bullet,G)$) is large.
\end{proposition}

\begin{proof}
    Suppose that $\Hom(\bullet,G)=\bigcup_{m<\omega}\mathcal H_m$ where each $\mathcal H_m$ is a tiny Borel set.
    For every $m<\omega$, Lemma~\ref{lemma:L0-claim3.5} produces an invariant Borel set $B_m\supseteq[\mathcal H_m]_{E_G}$ with a Borel 2-colouring $c_m$.
    Define $c\colon X\to2$ by letting $m(x):=\min\{m<\omega:x\in B_m\}$ and $c(x)=c_{m(x)}(x)$.
    Then $c$ is a Borel 2-colouring because each $B_m$ is invariant.
\end{proof}

Given that $L^c_{n+1}$ is obtained from two copies of $L^c_n$ joined by a path, it is convenient to have a means to refering to these distinct copies.
Keeping property \eqref{eq:L_n-condition-iii} in mind, for every $\varphi\in\Hom(L^c_{n+1},G)$ and $i<2$, define $\varphi^i\in\Hom(L^c_n,G)$ by $\varphi^i(u)=\varphi(u^\frown (i))$ for all $u\in V(L^c_n)$ and $\varphi^i(uv)=\varphi(u^\frown(i),v^\frown(i))$ for all $(u,v)\in E(L^c_n)$.
Thus, given $\mathcal L\subseteq\Hom(L^c_n,G)$, define $\mathcal L^{+c}$ as the set of all $\varphi\in\Hom(L^c_{n+1},G)$ for which $\varphi^i\in\mathcal L$ for both $i<2$.
Notice that if $\mathcal L$ is Borel, then so is $\mathcal L^{+c}$.

The length of the path added needs to be updated dynamically throughout the proof.
Formally, one adjusts the parameter $c$ by moving into a different branch of the Baire space when constructing the graphs $L^c_n$.
Notice that, for two $c,d\in\omega^\omega$ and for every $n<\omega$, the condition $d\restriction_n=c\restriction_n$ implies that $L^c_k=L^d_k$ for all $k<n$.

\begin{lemma}\label{lemma:Ln_basics}~ 
    \begin{enumerate}[label=(\alph*)]
        \item Suppose that $n>0$, that $c$ only takes odd values, that $k$ is a natural and $t\in2^{<n}$ is such that $(p_k)^\smallfrown t^\smallfrown(i)$, for $i<2$, are vertices of $L^c_n$ (see \eqref{eq:L_n-condition-iii}).
        Then, they are an odd distance apart in $L^c_n$.

        \item For any $n<\omega$, if $\mathcal L\subseteq\Hom(L^c_n,G)$ is a large Borel set, then for any $N<\omega$ there exists $d\in\omega^\omega$ such that $d\restriction_n=c\restriction_n$, $d(n)\geq N$ is odd, and $\mathcal L^{+d}$ is a nonempty subset of $\Hom(L^{d}_{n+1},G)$.

        \item For any $n<\omega$, if $\mathcal L\subseteq\Hom(L^c_n,G)$ is a large Borel set, then there exists $d\in\omega^\omega$ such that $d\restriction_n=c\restriction_n$ and $\mathcal L^{+d}$ is a large subset of $\Hom(L^{d}_{n+1},G)$.
    \end{enumerate}
\end{lemma}

\begin{proof}
    For (a), observe that both vertices come from the same vertex $(p_k)^\smallfrown t$ in $L^c_{|t|+1}$, and so their distance in $L^c_n$ is twice the distance from $(p_k)^\smallfrown t$ to $e_1^{|t|}$ in $L^c_{|t|+1}$ (since the two vertices are each in distinct copies of $L^c_{|t|+1}$ inside $L^c_n$ and the unique path between them passes through the joining path at stage $|t|$), plus the length of the joining path at that stage, which is $c(|t|)+2$.
    Since $c(|t|)$ is odd, this distance is odd.

    For (b), since $\mathcal L$ is large, it is not tiny.
    Hence, for every vertex $u\in V(L^c_n)$ and every natural $k$, $\Phi(\mathcal L(u),k)$ fails.
    In particular, fix the gluing vertex $u_0:=e_1^n\in V(L^c_n)$ (recall that $e_1^0=p_0$).
    Then for every $k$ there exist $\varphi_0,\varphi_1\in\mathcal L$ and an odd walk $(v_0,v_1,\dots,v_m)$ in $G$ with $v_0=\varphi_0(u_0)$, $v_m=\varphi_1(u_0)$, and $m>2k-1$.
    Choosing $k$ large enough, it is ensured that $m\geq N+2$ and $m$ is odd.
    Since $m$ is odd, $d(n):=m-2$ is also odd and satisfies $d(n)\geq N$.
    For each $1\leq j\leq m$, since $(v_{j-1},v_j)\in E(G)=\pi[E]$, pick $e_j\in E$ with $\pi(e_j)=(v_{j-1},v_j)$.
    Setting $d(x)=c(x)$ for all $x\neq n$ and noting $L^c_n=L^d_n$ (since $d\restriction_n=c\restriction_n$), define $\varphi\in\Hom(L_{n+1}^d,G)$ by $\varphi^0=\varphi_0$, $\varphi^1=\varphi_1$, $\varphi(p_i)=v_{i+1}$ for $0\leq i\leq m-2=d(n)$, and the edge assignments $\varphi(u_0^\frown (0),p_0)=e_1$, $\varphi(p_j,p_{j+1})=e_{j+1}$ for $0\leq j<m-2$, and $\varphi(p_{m-2},u_0^\frown (1))=e_m$.
    Then $\varphi\in\mathcal L^{+d}$.

    Next, prove (c) by contradiction.
    Recall that constructing $L^c_{n+1}$ from $L^c_n$ only requires knowing the value of $c(n)$.
    Suppose that for all $d$ that agree with $c$ except possibly in the $n$-th value, $\mathcal L^{+d}=\bigcup_m\mathcal F^d_m$ where each $\mathcal F^d_m$ is a tiny Borel subset of $\Hom(L^d_{n+1},G)$.
    For each $m$ and $d$, let $w^d_m\in V(L^d_{n+1})$ and $k^d_m<\omega$ be such that $\Phi(\mathcal F^d_m(w^d_m),k^d_m)$.
    
    Since $w^d_m\in V(L^d_{n+1})$, it is either a \emph{non-path vertex} of the form $u^\smallfrown(i)$ for some $u\in V(L^c_n)$ and $i<2$, or a \emph{path vertex} $p_j$ for some $0\leq j\leq d(n)$ in the joining path.
    I claim that in either case, there exists a vertex $u^d_m\in V(L^c_n)$ and $\ell^d_m<\omega$ such that $\Phi(\mathcal F^d_m({u^d_m}^\smallfrown(i^d_m)),\ell^d_m)$ for some $i^d_m<2$ where $u^d_m{}^\smallfrown(i^d_m)$ is a non-path vertex of $L^d_{n+1}$.

    Indeed, if $w^d_m=u^\smallfrown(i)$ is already non-path, take $u^d_m=u$, $i^d_m=i$, $\ell^d_m=k^d_m$.
    If $w^d_m=p_j$ is a path vertex of $L^d_{n+1}$, then the joining path connects $e_1^n{}^\smallfrown(0)$ to $e_1^n{}^\smallfrown(1)$ via $p_0,\dots,p_{d(n)}$.
    Since $p_j$ and $e_1^n{}^\smallfrown(0)$ are at distance $j+1$ in $L^d_{n+1}$, every homomorphism $\varphi\in\mathcal F^d_m$ maps them to vertices connected by a walk of lengh $j+1$ in $G$.
    Given any odd walk $(a_0,\dots,a_r)$ in $G$ with $a_0,a_r\in\mathcal F^d_m(e_1^n{}^\smallfrown(0))$, witnessed by $\varphi,\varphi'\in\mathcal F^d_m$, the walk $(\varphi(p_j),\dots,\varphi(e_1^n{}^\smallfrown(0))=a_0,\dots,a_r=\varphi'(e_1^n{}^\smallfrown(0)),\dots,\varphi'(p_j))$ has length $r+2(j+1)$, which has the same parity as $r$ (hence is also odd), and has endpoints in $\mathcal F^d_m(p_j)$.
    By $\Phi(\mathcal F^d_m(p_j),k^d_m)$, $r+2(j+1)\leq 2k^d_m-1$, so $r\leq 2(k^d_m-j-1)-1$.
    It follows that $\Phi(\mathcal F^d_m(e_1^n{}^\smallfrown(0)),\max\{0,k^d_m-j-1\})$.
    Thus, set $u^d_m=e_1^n$, $i^d_m=0$, $\ell^d_m=\max\{0,k^d_m-j-1\}$.
    
    \begin{claim}
        For each $m$ and $d$, there is a Borel set $B^d_m\supseteq\mathcal F^d_m({u^d_m}^\smallfrown(i^d_m))$ such that $\Phi(B^d_m,\ell^d_m)$.
    \end{claim}

    \begin{claimproof}
        Use the First Reflection Theorem (Lemma~\ref{lemma:first-reflection}).
        Fix $\ell=\ell^d_m$.
        It suffices to verify that $\Phi_\ell:=\{A\subseteq X:\Phi(A,\ell)\}$ is $\Pi^1_1$ on $\Sigma^1_1$.
        Let $Y$ be Polish and $A\subseteq Y\times X$ be $\Sigma^1_1$.
        Fix a Polish space $N$ and a continuous surjection $g\colon N\to A$.
        For each odd $r\geq 2\ell+1$, define
        \begin{align*}
            C_r:=\{(y,x_0,\dots,x_r,n_0,n_1,\hat e_1,\dots,\hat e_r)\in Y\times X^{r+1}\times N^2\times E^r:{}\\
            g(n_0)=(y,x_0)\land g(n_1)=(y,x_r)\land\textstyle\bigwedge_{j<r}\pi(\hat e_{j+1})=(x_j,x_{j+1})\}.
        \end{align*}
        Then $C_r$ is closed (by the continuity of $g$ and $\pi$), and $D_r:=\mathrm{proj}_Y(C_r)$ is $\Sigma^1_1$.
        Hence $D:=\bigcup\{D_r:r\text{ odd},\,r\geq 2\ell+1\}$ is $\Sigma^1_1$, and $A_{\Phi_\ell}=\{y\in Y:A_y\in\Phi_\ell\}=Y\setminus D$
        is $\Pi^1_1$.
        Since $\mathcal F^d_m({u^d_m}^\smallfrown(i^d_m))$ is analytic and satisfies $\Phi(\cdot,\ell)$, the First Reflection Theorem yields the desired Borel.
    \end{claimproof}

    Take $$\mathcal H^d_m=\{\varphi\in\mathcal L:\varphi(u^d_m)\in B^d_m\},$$ which is also a Borel subset of $\Hom(L^c_n,G)$.
    I claim that $\mathcal H^d_m$ is tiny.
    Indeed, $\mathcal H^d_m(u^d_m)=\{\varphi(u^d_m):\varphi\in\mathcal L\land\varphi(u^d_m)\in B^d_m\}\subseteq B^d_m$, so $\Phi(\mathcal H^d_m(u^d_m),\ell^d_m)$.
    
    Now, the set of pairs $(u^d_m,i^d_m)$ ranges over the finite set $V(L^c_n)\times2$ (which does not depend on $d$), and $(m,d)$ range over countably many values.
    Thus $\bigcup_{m,d}\mathcal H^d_m$ is a countable union of tiny sets, hence small.
    The set $\mathcal L_-:=\mathcal L\setminus\bigcup_{m,d}\mathcal H^d_m$ is therefore large (and Borel).
    By part (b), there exist $\varphi$ and $d$ such that $\varphi\in\mathcal L_-^{+d}\subseteq\mathcal L^{+d}$.
    Then $\varphi\in\mathcal F^d_m$ for some $m$.
    It follows that $\varphi^{i^d_m}(u^d_m)=\varphi({u^d_m}^\smallfrown(i^d_m))\in\mathcal F^d_m({u^d_m}^\smallfrown(i^d_m))\subseteq B^d_m$, hence $\varphi^{i^d_m}\in\mathcal H^d_m$.
    But $\varphi\in\mathcal L_-^{+d}$ implies $\varphi^{i^d_m}\in\mathcal L_-$, contradicting $\mathcal L_-\cap\mathcal H^d_m=\emptyset$.
\end{proof}

\subsection{Construction of a minimal graph of Borel chromatic number at least 3}\label{sec:L0-construction}

The goal of this subsection is to construct a family of graphs $\mathbb L_c$ indexed by $c\in\omega^\omega$.

Fix $c$ and consider $X_c$ as the set of all tuples $(m,k,x)\in\mathbb N\times\mathbb N\times2^\mathbb N$ such that either $m=0$ and $k=0$, or $m\geq1$ and $k\leq c(m-1)$; this is a closed subspace of the product space, and hence Polish.
Define, for each $(m,k,x)$ with $m\leq n$, $\pi_n(m,k,x)$ as the vertex of $L^c_n$ determined by these parameters.
Formally, $\pi_n(m,k,x):=(p_k)^\smallfrown x\restriction_{n-m}$.
For example, $\pi_3(1,2,0110\cdots)=(p_2,0)$ can be seen in Figure~\ref{fig:L_n} labeled as $(2,0)$.

Finally, $\mathbb L_c$ is the graph on $X_c$ where two $(n_i,k_i,x_i)$, for $i<2$, are adjacent if and only if the $\pi_n(n_i,k_i,x_i)$ are adjacent in all $L^c_n$ with $n\geq\max\{n_0,n_1\}$.
Some basic properties of $\mathbb L_c$ follow.

\begin{lemma}\label{lemma:L0_basics}~ 
    \begin{enumerate}[label=(\alph*)]
        \item Two vertices $(n_i,k_i,x_i)\in X_c$, for $i<2$, are in the same connected component of $\mathbb L_c$ if and only if there are $t_i\in2^{<\omega}$ and $x\in2^\omega$ such that $|t_0|-|t_1|=n_1-n_0$ and $x_i=t_i^\smallfrown x$.
        
        \item $\mathbb L_c$ has no cycles and all vertices of $\mathbb L_c$ have degree 2 except for the vertex $(0,0,\{0\}^\omega)$, which has degree 1.
        
        \item If, additionally, $c$ takes only odd values, then $\chi_B(\mathbb L_c)=3$.
    \end{enumerate}
\end{lemma}

\begin{proof}
    For (a), let $(n_i,k_i,x_i)\in X_c$ for $i<2$.
    If they are in the same connected component, then for all large enough $n$, $\pi_n(n_0,k_0,x_0)$ and $\pi_n(n_1,k_1,x_1)$ lie in the same copy of $L^c_{n-1}$ inside $L^c_n$.
    By the construction of $L^c_n$, two vertices $(p_{k_0})^\smallfrown x_0\restriction_{n-n_0}$ and $(p_{k_1})^\smallfrown x_1\restriction_{n-n_1}$ are in the same copy of $L^c_{n-1}$ in $L^c_n$ if and only if their last binary digits agree, that is, $x_0(n-n_0-1)=x_1(n-n_1-1)$ for all large $n$.
    Writing $x_0=t_0^\smallfrown x$ and $x_1=t_1^\smallfrown x$ with $|t_0|=n_1-n_0+|t_1|$ (assuming $n_0\leq n_1$), one obtains the stated condition.
    Conversely, if $x_i=t_i^\smallfrown x$ with $|t_0|-|t_1|=n_1-n_0$, then for all large $n$, $\pi_n(n_0,k_0,x_0)=(p_{k_0})^\smallfrown t_0^\smallfrown x\restriction_r$ and $\pi_n(n_1,k_1,x_1)=(p_{k_1})^\smallfrown t_1^\smallfrown x\restriction_r$ for some $r$, so they share a common tail and lie in the same connected component of $L^c_n$, hence of $\mathbb L_c$.

    For (b), every vertex of $L^c_n$ has degree at most $2$, so the same holds for $\mathbb L_c$.
    It follows from acyclicity of the finite paths $L^c_n$ that $\mathbb L_c$ is acyclic.
    The vertex $(0,0,0^\omega)$ maps under $\pi_n$ to $e_0^n=(p_0)^\smallfrown 0^n$, which is an endpoint of $L^c_n$, hence has degree $1$ in $L^c_n$ for all $n\geq1$.
    Thus $(0,0,0^\omega)$ has degree $1$ in $\mathbb L_c$.
    For any other vertex $(m,k,x)\in X_c$, the image $\pi_n(m,k,x)$ is an interior vertex of $L^c_n$ for all sufficiently large $n$, so it has degree exactly $2$.

    (c) That $\chi_B(\mathbb L_c)\leq3$ follows from Proposition~\ref{prop:brooks} and part (b).
    Argue the other inequality by contradiction.
    Suppose that there is a Borel partition of $X_c$ into two independent sets $X_c=I_0\cup I_1$.
    Since $X_c$ is Polish, the Baire category theorem guarantees that some $I_\varepsilon$ is non-meager, hence comeager in a basic open set $U=\{(m,k,x)\in X_c:x\restriction_r=t\}$ for some $m$, $k$, $r$, and $t\in2^r$.
    In particular, for a comeager set of $x\in2^\mathbb N$ extending $t$, $(m,k,x)\in I_\varepsilon$.
    Since $(m,k,t^\smallfrown(0)^\smallfrown y)$ and $(m,k,t^\smallfrown(1)^\smallfrown y)$ are in the same connected component (by part (a)) for any $y\in2^\mathbb N$, and $I_\varepsilon$ is comeager in $U$, one can find $y$ such that both $(m,k,t^\smallfrown(0)^\smallfrown y)$ and $(m,k,t^\smallfrown(1)^\smallfrown y)$ lie in $I_\varepsilon$.
    But by Lemma~\ref{lemma:Ln_basics}(a), these two vertices are an odd distance apart in $\mathbb L_c$ (as $(p_k)^\smallfrown t^\smallfrown(0)$ and $(p_k)^\smallfrown t^\smallfrown(1)$ are an odd distance apart in every $L^c_n$ for large enough $n$), contradicting $I_\varepsilon$ being independent.
\end{proof}

The main result is then split into two parts.

\begin{theorem}\label{thm:L0-hom-equivalent-main}
    For any analytic graph $G$, either $\chi_B(G)\leq 2$ or there is a $c$ such that $\mathbb L_c\to_cG$.
\end{theorem}

Moreover, $c$ can be chosen to be unbounded and take only odd values.

\begin{theorem}\label{thm:L0-hom-equivalent}
    Let $c,d\in\omega^\omega$ be unbounded and only taking odd values, then $\mathbb L_c$ and $\mathbb L_d$ are continuously homomorphically equivalent.
\end{theorem}

The proof of Theorem~\ref{thm:L0-hom-equivalent} is purely combinatorial: one constructs homomorphisms between $\mathbb L_c$ and $\mathbb L_d$ by iteratively extending partial maps along finite path components, using the ``large gap property'' guaranteed by unboundedness.
The argument does not depend on the method used to establish Theorem~\ref{thm:L0-hom-equivalent-main}; see \cite[Proposition~4.2 and Claim~4.1]{L0paper} for details.

Combining Theorems~\ref{thm:L0-hom-equivalent-main} and \ref{thm:L0-hom-equivalent}, and fixing any unbounded odd $c_0$ (e.g., $c_0(0)=1$ and $c_0(n)=2n-1$ for $n>0$)\footnote{It is perhaps worth noting that I picked a \emph{different} parameter for Figure~\ref{fig:L_n} purely for aesthetic purposes.}, the $L_0$ dichotomy of \cite{L0paper} is recovered:

\begin{corollary}
    Let $\mathbb L_0:=\mathbb L_{c_0}$.
    For any analytic graph $G$ on a Polish space, exactly one of the following holds:
    \begin{enumerate}[label=(\alph*)]
        \item $\chi_B(G)\leq 2$;
        \item $\mathbb L_0\to_cG$.
    \end{enumerate}
\end{corollary}

\subsection{Proof of Theorem~\ref{thm:L0-hom-equivalent-main}}

By Lemma~\ref{lemma:L0_basics}(c) and Fact~\ref{fact:chrom_transitive}, it is clear that both conditions cannot hold simultaneously.
Now suppose that $G$ is an analytic graph of Borel chromatic number at least three.
The theorem is proved by constructing a valid $c$ and a continuous homomorphism witnessing $\mathbb L_c\to_cG$.

By Proposition~\ref{prop:small}, $X$ is large.
Fix compatible complete metrics $\rho$ on $X$ and $\delta$ on $E$.
Since $L^c_0=\bullet$, $\Hom(L^c_0,G)$ may be identified with $X$ and set $\mathcal L_0:=X$, which is large.
Recursively construct a function $c\in\omega^\omega$ (taking only odd values) and a sequence of large Borel sets $\mathcal L_n\subseteq\Hom(L^c_n, G)$ such that for all $n$,

\begin{enumerate}[label=(\arabic*)]
    \item $\mathcal L_{n+1}\subseteq\mathcal L_n^{+c}$;
    \item for all $u\in V(L^c_n)$, $\mathrm{diam}_\rho(\overline{\mathcal L_n(u)})< 2^{-n}$; and
    \item for all $(u,v)\in E(L^c_n)$, $\mathrm{diam}_\delta(\overline{\mathcal L_n(u,v)})< 2^{-n}$.
\end{enumerate}

Given $\mathcal L_n$ large, apply Lemma~\ref{lemma:Ln_basics}(c) to obtain $d\in\omega^\omega$ with $d\restriction_n=c\restriction_n$ and $\mathcal L_n^{+d}$ large; set $c(n):=d(n)$.
To arrange conditions (2) and (3) for $n+1$: since $X$ and $E$ can each be covered by countably many closed sets of $\rho$-diameter (respectively $\delta$-diameter) less than $2^{-(n+1)}$, and small sets form a $\sigma$-ideal, one may intersect with the preimages of these covers and discard the resulting small pieces, passing to a large Borel subset $\mathcal L_{n+1}\subseteq\mathcal L_n^{+c}$ satisfying (2) and (3) at stage $n+1$.
Moreover, by choosing $c(n)$ sufficiently large using part (b) at each stage, $c$ can be made unbounded.

By (1), for all $n\geq m$, $$\mathcal L_{n+1}(\pi_{n+1}(m,k,x))\subseteq\mathcal L_n(\pi_n(m,k,x));$$ indeed, given that $\pi_{n+1}(m,k,x)=\pi_n(m,k,x)^\frown(x(n-m))$, each $\varphi\in\mathcal L_{n+1}$ satisfies $\varphi^{x(n-m)}\in\mathcal L_n$ and $$\varphi^{x(n-m)}(\pi_n(m,k,x))=\varphi(\pi_{n+1}(m,k,x)).$$
By the same reasoning, if $\pi_n(m_0,k_0,x_0)$ and $\pi_n(m_1,k_1,x_1)$ are adjacent in $L^c_n$ for all large $n$, then \begin{align*}
    \mathcal L_{n+1}(\pi_{n+1}(m_0,k_0,x_0),\pi_{n+1}(m_1,k_1,x_1))\\
    \subseteq\mathcal L_n(\pi_n(m_0,k_0,x_0),\pi_n(m_1,k_1,x_1)).
\end{align*}

Thus, the sets $\overline{\mathcal L_n(\pi_n(m,k,x))}$ are closed and nested, with diameters tending to zero (see (2) from before).
Since $\rho$ is a complete metric, Cantor's intersection theorem applies.
For $(m,k,x)\in X_c$, let $f(m,k,x)$ be the unique point in $$\bigcap_{n=m}^\infty\overline{\mathcal L_n(\pi_n(m,k,x))}=\bigcap_{n=0}^\infty\overline{\mathcal L_{n+m}((p_k)^\frown x\restriction_n)}.$$
The map $f\colon X_c\to X$ is continuous.

It remains to check that it is a homomorphism from $\mathbb L_c$ to $G$.
Suppose that $(m_i,k_i,x_i)_{i<2}$ are adjacent in $\mathbb L_c$, that is, for all $n\geq n_0:=\max\{m_0,m_1\}$, $\pi_n(m_i,k_i,x_i)$ are adjacent in $L_n^c$.
Similar to the preceding paragraph, the fact that $\delta$ is a complete metric implies that there is a unique edge $e\in E$ inside $$\bigcap_{n\geq n_0}\overline{\mathcal L_n(\pi_n(m_0,k_0,x_0),\pi_n(m_1,k_1,x_1))}.$$
By the continuity of $\pi$, $(f(m_0,k_0,x_0),f(m_1,k_1,x_1))=\pi(e)\in E(G)$, in other words, the images form an edge of $G$, and this concludes the proof.

The hereby presented argument shows that the core of the $L_0$ dichotomy, namely the construction of a continuous homomorphism from an $\mathbb{L}_c$ graph into any analytic graph of Borel chromatic number at least three, admits a proof that closely parallels Bernshteyn's argument for the $G_0$ dichotomy, with the challenge of adding a dynamically-changing parameter to the construction.

The main new ingredient is the replacement of the independence-based $\sigma$ (sets whose vertex projections are independent) with the odd-walk-bounded $\sigma$-ideal defined by the $\Phi$ property.
It remains an interesting question whether analogous adaptations can yield simplified proofs of other dichotomy results in Borel combinatorics, or even entirely new ones.